\begin{example}[A field]\label{ex:vi28-01}
Every field is integrally closed in itself, so $\Spec K$ is normal.
\end{example}

\begin{example}[The integers]\label{ex:vi28-02}
$\mathbb Z$ is a UFD and hence integrally closed in $\mathbb Q$; therefore $\Spec\mathbb Z$ is normal.
\end{example}

\begin{example}[Affine space]\label{ex:vi28-03}
$k[x_1,\dots,x_n]$ is a UFD, hence normal, so $\mathbb A^n_k$ is normal.
\end{example}

\begin{example}[Principal opens]\label{ex:vi28-04}
If $A$ is normal and $f\in A$, then $A_f$ is normal; every distinguished open of a normal affine scheme is normal.
\end{example}

\begin{example}[The cusp]\label{ex:vi28-05}
$A=k[t^2,t^3]$ is integral but not normal because $t\in\Frac(A)$ is integral over $A$ and $t\notin A$.
\end{example}

\begin{example}[Normalization of the cusp]\label{ex:vi28-06}
The integral closure of $k[t^2,t^3]$ in $k(t)$ is $k[t]$, giving $\mathbb A^1_k\to V(y^2-x^3)$.
\end{example}

\begin{example}[A quadratic order]\label{ex:vi28-07}
The integral closure of $\mathbb Z[\sqrt5]$ in $\mathbb Q(\sqrt5)$ is $\mathbb Z[(1+\sqrt5)/2]$.
\end{example}

\begin{example}[A DVR]\label{ex:vi28-08}
Every discrete valuation ring is integrally closed, hence its spectrum is normal.
\end{example}

\begin{example}[Already normal]\label{ex:vi28-09}
If $A=\overline A$, then $\Spec\overline A\to\Spec A$ is the identity up to canonical isomorphism.
\end{example}

\begin{example}[Localization of a cusp]\label{ex:vi28-10}
Away from the singular point, the cusp is already normal; normalization is an isomorphism over that open set.
\end{example}

\begin{example}[Function field unchanged]\label{ex:vi28-11}
For any affine normalization $A\subseteq\overline A\subseteq\Frac(A)$, one has $\Frac(\overline A)=\Frac(A)$.
\end{example}

\begin{example}[Surjectivity]\label{ex:vi28-12}
Because $\overline A$ is integral over $A$, lying over makes $\Spec\overline A\to\Spec A$ surjective.
\end{example}

\begin{example}[Dominance]\label{ex:vi28-13}
The inclusion $A\hookrightarrow\overline A$ has zero kernel, so the normalization morphism is dominant.
\end{example}

\begin{example}[Nodal curve]\label{ex:vi28-14}
The normalization of $y^2=x^2(x+1)$ separates the two branches through the node; two points of the normal curve lie over the singular point.
\end{example}

\begin{example}[Normal curve chart]\label{ex:vi28-15}
A nonsingular affine line chart $k[t]$ is already normal, so normalization does nothing there.
\end{example}

\begin{example}[Integral but not normal]\label{ex:vi28-16}
Integrality only asks for a domain; $k[t^2,t^3]$ shows that a domain can fail integral closedness.
\end{example}

\begin{example}[Finite normalization]\label{ex:vi28-17}
If $A$ is a finitely generated $k$-domain, then $\overline A$ is finite over $A$, so normalization is a finite morphism.
\end{example}

\begin{example}[Reduction is different]\label{ex:vi28-18}
$\Spec k[\varepsilon]/(\varepsilon^2)$ has a reduction but is not integral, so the integral-scheme normalization construction is not the same operation.
\end{example}

\begin{example}[Normality is local]\label{ex:vi28-19}
If an integral scheme has an affine cover by spectra of normal domains, then all of its local rings are integrally closed.
\end{example}

\begin{example}[Normal open subsets]\label{ex:vi28-20}
Every nonempty open subscheme of a normal integral scheme is again normal and has the same function field.
\end{example}
